\documentclass[11pt,letterpaper]{article}
% Reusable style. Build each .tex from its own directory using scripts/build.sh.
\usepackage[T1]{fontenc}
\usepackage[utf8]{inputenc}
\usepackage{lmodern}
\usepackage[margin=0.86in,headheight=15pt]{geometry}
\usepackage{amsmath,amssymb,amsthm,mathtools}
\usepackage{booktabs,longtable,tabularx,array,multicol}
\usepackage[table]{xcolor}
\usepackage{microtype}
\usepackage{enumitem}
\usepackage{fancyhdr}
\usepackage{titlesec}
\usepackage{hyperref}
\usepackage{listings}
\usepackage[most]{tcolorbox}
\usepackage{needspace}
\definecolor{aeNavy}{HTML}{172D43}
\definecolor{aeBlue}{HTML}{176B90}
\definecolor{aeTeal}{HTML}{148477}
\definecolor{aePale}{HTML}{F2F7F9}
\definecolor{aeGray}{HTML}{5C6771}
\hypersetup{colorlinks=true,linkcolor=aeBlue,urlcolor=aeBlue,pdfauthor={Algorithm Engineering Documentation Suite},pdfsubject={Independent educational materials based on the supplied TAOCP volume topics}}
\titleformat{\section}{\Large\bfseries\color{aeNavy}}{\thesection}{0.6em}{}
\titleformat{\subsection}{\large\bfseries\color{aeBlue}}{\thesubsection}{0.6em}{}
\titleformat{\subsubsection}{\normalsize\bfseries\color{aeNavy}}{\thesubsubsection}{0.6em}{}
\setlist[itemize]{leftmargin=*,topsep=3pt,itemsep=1pt}
\setlist[enumerate]{leftmargin=*,topsep=3pt,itemsep=3pt}
\setlength{\parindent}{0pt}\setlength{\parskip}{5pt}
\pagestyle{fancy}\fancyhf{}\fancyhead[L]{\footnotesize\color{aeGray}ALGORITHM ENGINEERING}\fancyhead[R]{\footnotesize\color{aeGray}DOCUMENTATION SUITE}\fancyfoot[L]{\footnotesize\color{aeGray}Independent study and engineering practice}\fancyfoot[R]{\footnotesize\thepage}
\renewcommand{\headrulewidth}{0.35pt}
\newtcolorbox{engineeringnote}[1][]{breakable,colback=aePale,colframe=aeBlue,boxrule=0.5pt,arc=1mm,fonttitle=\bfseries,title={Engineering note},#1}
\newtcolorbox{evidencebox}[1][]{breakable,colback=white,colframe=aeTeal,boxrule=0.6pt,arc=1mm,fonttitle=\bfseries,title={Evidence and deliverables},#1}
\lstdefinestyle{algo}{basicstyle=\ttfamily\footnotesize,breaklines=true,frame=single,rulecolor=\color{aeGray!40},backgroundcolor=\color{aePale},showstringspaces=false,tabsize=2,columns=fullflexible,keepspaces=true,keywordstyle=\color{aeBlue}\bfseries,commentstyle=\color{aeTeal}}
\lstset{style=algo}
\newcommand{\DocTitle}[3]{%
\hypersetup{pageanchor=false}\begin{titlepage}\vspace*{1.2cm}
{\color{aeTeal}\rule{\linewidth}{2.3pt}}\par\vspace{1.15cm}
{\fontsize{24}{30}\selectfont\bfseries\color{aeNavy}\raggedright #1\par}\vspace{0.45cm}
{\Large\color{aeBlue}#2\par}\vspace{1.3cm}
{\large\bfseries #3\par}\vspace{0.35cm}
{\color{aeGray}Edition 1.0 \quad | \quad 29 October 2026\par}\vfill
{\color{aeGray}\small Original educational material aligned to the chapter and section topics in the five supplied volumes of Donald E. Knuth's \emph{The Art of Computer Programming}. This document is not an authorized reproduction or substitute for the original books.}\par\vspace{0.4cm}
{\color{aeTeal}\rule{\linewidth}{0.75pt}}
\end{titlepage}\hypersetup{pageanchor=true}\pagenumbering{roman}}
\newcommand{\SourceNote}[1]{\par\smallskip{\footnotesize\color{aeGray}\textbf{Source locator:} #1}\par\smallskip}
\newcommand{\Outcome}[1]{\par\textbf{Competency target.} #1\par}
\newcommand{\Checkpoint}[1]{\par\textbf{Checkpoint.} #1\par}
% Intentionally simple printable writing lines: call \AnswerLines{number}.
\newcount\alcount
\newcommand{\AnswerLines}[1]{\par\alcount=0\loop\ifnum\alcount<#1\par\noindent\rule{\linewidth}{0.3pt}\vspace{0.33cm}\advance\alcount by1\repeat\par}
\newcommand{\FormField}[2][0.62\linewidth]{\noindent\textbf{#2}\hspace{0.65em}\rule{#1}{0.4pt}\par\vspace{0.18cm}}
\newcommand{\ReviewTable}{\begin{tabularx}{\linewidth}{@{}p{0.23\linewidth}X p{0.2\linewidth}@{}}\toprule Criterion & Recorded evidence & Rating (0--4)\\\midrule Correctness & & \\[1em] Proof quality & & \\[1em] Complexity & & \\[1em] Testing & & \\[1em] Reproducibility & & \\[1em] Communication & & \\[1em]\bottomrule\end{tabularx}}

\begin{document}
\DocTitle{Volume 4B Study Guide}{Chapter 7 continuation: backtracking, exact cover and satisfiability}{Chapter- and section-aligned study guide}
\tableofcontents
\clearpage
\pagenumbering{arabic}
\section{Study protocol}

For each subsection, read the specified part of the original volume, keep a mathematical notebook, produce an original implementation or design sketch, and evaluate the acceptance gate. The investigations below are new exercises, not copies of the book's problems or solutions. Annotate any starred or mathematically advanced source sections as optional on the first pass; revisit them after the prerequisite mathematics is in place.
\begin{evidencebox}Recommended artifacts: study notes, problem specification, derivation, runnable code or model, tests, performance evidence and reflective comparison to original assumptions.\end{evidencebox}

\section{\S 7.2.2: Backtrack programming}
\SourceNote{Knuth, Volume 4B, \S 7.2.2.}
\Outcome{Build reversible state updates, choice ordering, pruning and explicit solution counters.}
\subsection{Independent investigation}
Solve n-queens by backtracking, instrumenting nodes and failed branches for n=4 to 10.
\subsection{Analysis and verification gate}
Cross-check known small counts via independent brute force and explain branching growth.
\subsection{Study questions}
\begin{enumerate}
\item Which representations, input assumptions and cost metrics control this technique's applicability?
\item Which invariant or combinatorial identity establishes correctness, and which input is its strongest counterexample candidate?
\item What result would cause the implementation to fail review despite passing a small collection of examples?
\end{enumerate}
\Checkpoint{Preserve a short explanation, the relevant test evidence, and one unresolved question before proceeding.}

\section{\S 7.2.2.1: Dancing links and exact cover}
\SourceNote{Knuth, Volume 4B, \S 7.2.2.1.}
\Outcome{Model primary items, options, reversible cover/uncover, secondary items and colored extensions.}
\subsection{Independent investigation}
Build a simple set-based exact cover solver first, then compare with a reversible linked representation.
\subsection{Analysis and verification gate}
Verify that every reported solution covers each primary item once; test impossible and multiple-solution instances.
\subsection{Study questions}
\begin{enumerate}
\item Which representations, input assumptions and cost metrics control this technique's applicability?
\item Which invariant or combinatorial identity establishes correctness, and which input is its strongest counterexample candidate?
\item What result would cause the implementation to fail review despite passing a small collection of examples?
\end{enumerate}
\Checkpoint{Preserve a short explanation, the relevant test evidence, and one unresolved question before proceeding.}

\section{\S 7.2.2.2: Satisfiability}
\SourceNote{Knuth, Volume 4B, \S 7.2.2.2.}
\Outcome{Model CNF, literal semantics, unit propagation, branching, and distinctions among chronological backtracking and learning.}
\subsection{Independent investigation}
Implement a small DPLL solver, encode graph 2-coloring and compare SAT answers against exhaustive assignments.
\subsection{Analysis and verification gate}
For SAT outputs validate the returned witness; for UNSAT small formulas cross-check exhaustive search.
\subsection{Study questions}
\begin{enumerate}
\item Which representations, input assumptions and cost metrics control this technique's applicability?
\item Which invariant or combinatorial identity establishes correctness, and which input is its strongest counterexample candidate?
\item What result would cause the implementation to fail review despite passing a small collection of examples?
\end{enumerate}
\Checkpoint{Preserve a short explanation, the relevant test evidence, and one unresolved question before proceeding.}

\section{Volume synthesis and mastery review}
\begin{enumerate}
\item Draw a dependency graph connecting at least four major ideas from this volume.
\item Select two alternative implementations for one workload and compare their preconditions, proven cost, memory use, measured performance and failure cases.
\item Present a five-minute technical defense of the most surprising experimental result. A claim that differs from the mathematical model requires a measured or documented explanation.
\end{enumerate}
\subsection{Completion checklist}
\begin{itemize}\item All module notes include a source locator and an independently stated learning outcome.
\item Every implementation has reproducible tests and edge cases.
\item Every complexity claim identifies its cost model and underlying assumptions.
\item A peer can independently reproduce one chosen experiment using its recorded configuration.
\end{itemize}
\end{document}
